% Part: computability
% Chapter: recursive-functions
% Section: introduction

\documentclass[../../../include/open-logic-section]{subfiles}

\begin{document}

\olfileid{cmp}{rec}{int}
\olsection{परिचय}

संगणनक्षमतेची गणितीय उपपत्ती विकसित करण्यासाठी सर्वप्रथम
संगणनक्षमतेचे एक \emph{प्रतिमान} विकसित करावे लागते. आज आपण
संगणनक्षमता ही संगणक करतात त्या प्रकारची गोष्ट आहे असे मानतो, आणि संगणक
चिन्हांच्या साहाय्याने काम करतात. परंतु संगणनक्षमतेच्या उपपत्तींच्या
विकासाच्या आरंभी संगणनाचे आदर्श उदाहरण \emph{संख्यात्मक} संगणन होते.
गणितज्ञांना संख्या-सैद्धांतिक फलनांमध्ये, म्हणजे संगणित करता येणाऱ्या
$f\colon \Nat^n \to \Nat$ या फलनांमध्ये, नेहमीच रस होता. म्हणून
संगणनक्षमतेच्या उपपत्तीच्या आरंभी अशाच फलनांचा अभ्यास झाला, यात आश्चर्य
नाही. बेरीज, गुणाकार, (नैसर्गिक संख्यांची) घातांक क्रिया यांसारख्या
संगणनक्षम संख्यात्मक फलनांच्या सर्वाधिक परिचित उदाहरणांत एक रोचक समान
वैशिष्ट्य आहे: त्यांची \emph{पुनरावर्ती रीतीने} व्याख्या करता येते.
त्यामुळे पुनरावर्ती व्याख्यांच्या आधारे \emph{संगणनक्षम फलनाची}
सर्वसाधारण व्याख्या करण्याचा प्रयत्न अगदी स्वाभाविक आहे.
संख्या-सैद्धांतिक फलनांची पुनरावर्ती रीतीने व्याख्या करण्याच्या अनेक
संभाव्य मार्गांपैकी येथे एका विशेषतः सोप्या व्याख्या-आकृतिबंधाला
मध्यवर्ती स्थान मिळते: तथाकथित \emph{आदिम पुनरावर्तन}.

संगणनक्षम फलनांबरोबरच संगणनक्षम संच आणि संबंध यांतही आपल्याला रस असू
शकतो. दिलेली संख्या संचाचा !!a{element} आहे की नाही याचे उत्तर आपण
संगणित करू शकत असू, तर तो संच संगणनक्षम असतो; आणि दिलेला क्रमित
घटकसमूह $\tuple{n_1, \dots, n_k}$ हा संबंधाचा !!a{element} आहे
की नाही हे आपण संगणित करू शकत असू तेव्हा आणि केवळ तेव्हाच तो संबंध
संगणनक्षम असतो. संचाचे किंवा संबंधाचे \emph{जातिबोधक फल} विचारात
घेतल्यास संगणनक्षम संच व संबंध यांची चर्चा संगणनक्षम फलनांच्या चर्चेत
समाविष्ट करता येते. त्यामुळे आदिम पुनरावर्ती संबंधांचीही व्याख्या करता
येते; उदाहरणार्थ, ``$m$ ला $n$ ने निःशेष भाग जातो'' हा संबंध आदिम
पुनरावर्ती आहे.

तथापि, आदिम पुनरावर्ती फलने---केवळ आदिम पुनरावर्तन वापरून परिभाषित करता
येणारी फलने---हीच एकमेव संगणनक्षम संख्या-सैद्धांतिक फलने नाहीत. आदिम
पुनरावर्तनाची अनेक सामान्यीकरणे विचारात घेण्यात आली आहेत; परंतु फलने
संगणित करण्याची सर्वाधिक शक्तिशाली आणि व्यापक मान्यता असलेली अतिरिक्त
पद्धत म्हणजे अपरिबद्ध शोध. यातून \emph{आंशिक पुनरावर्ती फलनांची}
व्याख्या मिळते, आणि त्यासंबंधित व्याख्येतून \emph{सामान्य पुनरावर्ती
  फलनांची} व्याख्या मिळते. सामान्य पुनरावर्ती फलने संगणनक्षम आणि
सर्वत्र परिभाषित असतात; आणि ही व्याख्या आंशिक पुनरावर्ती फलनांपैकी
सर्वत्र परिभाषित असणारी नेमकी फलने लक्षणित करते. पुनरावर्ती फलने
संगणनाच्या प्रत्येक इतर प्रतिमानाचे (ट्यूरिंग यंत्रे, लॅम्डा कलन इत्यादींचे)
अनुकरण करू शकतात आणि म्हणून ती संगणनाच्या अनेक मान्य प्रतिमानांपैकी एक
प्रतिमान ठरतात.


\end{document}
